\documentclass[11pt,a4paper]{article}
\usepackage{fontspec}\usepackage[french]{babel}
\usepackage[margin=1.8cm]{geometry}\usepackage{tabularx,colortbl,xcolor,array}
\setmainfont{PlusJakartaSans-Medium.ttf}[Path=fonts/,BoldFont=PlusJakartaSans-Bold.ttf]
\usepackage{xeCJK}\setCJKmainfont{WenQuanYiZenHei-subset.ttf}[Path=fonts/]
\pagestyle{empty}\setlength{\parindent}{0pt}
\begin{document}
{\LARGE\bfseries Chinois — Terminale A\par}\smallskip
{\large Répartition des 40 leçons par chapitre et par trimestre\par}\medskip
Cours « signés OnBuch+ », conformes à la fiche de progression harmonisée du MINESEC (année scolaire 2026-2027). Chaque leçon est un PDF autonome (3486 pages au total pour les 40 leçons).\par\medskip
\textbf{Contenu type d'une leçon :} page de garde, sommaire, objectifs, prérequis, sections de cours (textes en chinois simplifié avec pinyin et traduction, vocabulaire en tableaux, structures de grammaire, clés et ordre des traits, exemples traduits), activité d'intégration, méthodes et exercices résolus, exercices (connaissances, entraînement, préparation à l'examen), corrigés, fiche bilan.\par\medskip
\textbf{Les 5 modules :} Module 1 : Vie familiale et intégration sociale (leçons 1--8) ; Module 2 : Environnement et santé (leçons 9--16) ; Module 3 : Citoyenneté et ouverture au monde (leçons 17--24) ; Module 4 : Activités économiques et monde du travail (leçons 25--33) ; Module 5 : Mass médias et télécommunication (leçons 34--40).\par\bigskip
\section*{Module 1 : Vie familiale et intégration sociale}
\textit{8 leçons, leçons 1 à 8, 694 pages.}\par\medskip
\begin{tabularx}{\linewidth}{|c|X|l|l|c|}\hline
\rowcolor{gray!20} N° & Titre de la leçon & Trimestre & Période & Pages \\ \hline
1 & Leçon 1 — Vocabulaire et compréhension de texte : situation de la femme (égalité homme-femme) \newline {\small\emph{Type : vocabulaire — fichier : ZH01-vocabulaire-situation-de-la-femme-egalite-homme-fe.pdf}} & 1er trimestre & 07/09 – 02/10/2026 & 67 \\ \hline
2 & Leçon 2 — Grammaire : le comparatif;……跟/和/与......一样/相同;……跟/和/与......不一样/不同。 \newline {\small\emph{Type : grammaire — fichier : ZH02-grammaire-le-comparatif.pdf}} & 1er trimestre & 07/09 – 02/10/2026 & 82 \\ \hline
3 & Leçon 3 — Expression écrite : situation de la femme \newline {\small\emph{Type : ecrit — fichier : ZH03-ecrit-situation-de-la-femme.pdf}} & 1er trimestre & 07/09 – 02/10/2026 & 99 \\ \hline
4 & Leçon 4 — Éléments socioculturels : processus d’hérédité au Cameroun et en Chine \newline {\small\emph{Type : socio — fichier : ZH04-socio-processus-dheredite-au-cameroun-et-en-chine.pdf}} & 1er trimestre & 07/09 – 02/10/2026 & 89 \\ \hline
5 & Leçon 5 — Vocabulaire et compréhension de texte : héritage \newline {\small\emph{Type : vocabulaire — fichier : ZH05-vocabulaire-heritage.pdf}} & 1er trimestre & 07/09 – 02/10/2026 & 82 \\ \hline
6 & Leçon 6 — Grammaire : ……才/ 就...... Et 还 ; 又 ; 再 \newline {\small\emph{Type : grammaire — fichier : ZH06-grammaire-et.pdf}} & 1er trimestre & 07/09 – 02/10/2026 & 76 \\ \hline
7 & Leçon 7 — Expression écrite : l'héritage \newline {\small\emph{Type : ecrit — fichier : ZH07-ecrit-l-heritage.pdf}} & 1er trimestre & 07/09 – 02/10/2026 & 97 \\ \hline
8 & Leçon 8 — Éléments socioculturels : processus d'hérédité au Cameroun et en Chine (suite : comparaison et exposé) \newline {\small\emph{Type : socio — fichier : ZH08-socio-processus-dheredite-au-cameroun-et-en-chine.pdf}} & 1er trimestre & 07/09 – 02/10/2026 & 102 \\ \hline
\end{tabularx}\bigskip
\section*{Module 2 : Environnement et santé}
\textit{8 leçons, leçons 9 à 16, 709 pages.}\par\medskip
\begin{tabularx}{\linewidth}{|c|X|l|l|c|}\hline
\rowcolor{gray!20} N° & Titre de la leçon & Trimestre & Période & Pages \\ \hline
9 & Leçon 9 — Vocabulaire et compréhension de texte : processus du recyclage \newline {\small\emph{Type : vocabulaire — fichier : ZH09-vocabulaire-processus-du-recyclage.pdf}} & 1er trimestre & 05/10 – 06/11/2026 & 108 \\ \hline
10 & Leçon 10 — Grammaire : Phrases interrogatives alternatives …还是… ; 首先…其次… \newline {\small\emph{Type : grammaire — fichier : ZH10-grammaire-phrases-interrogatives-alternatives.pdf}} & 1er trimestre & 05/10 – 06/11/2026 & 83 \\ \hline
11 & Leçon 11 — Expression écrite : le recyclage \newline {\small\emph{Type : ecrit — fichier : ZH11-ecrit-le-recyclage.pdf}} & 1er trimestre & 05/10 – 06/11/2026 & 90 \\ \hline
12 & Leçon 12 — Vocabulaire et compréhension de texte : paludisme \newline {\small\emph{Type : vocabulaire — fichier : ZH12-vocabulaire-paludisme.pdf}} & 1er trimestre & 05/10 – 06/11/2026 & 81 \\ \hline
13 & Leçon 13 — Grammaire : Interrogation avec 怎么了et怎么 ; 不但……, 而且…… ; 因为……, 所以…… \newline {\small\emph{Type : grammaire — fichier : ZH13-grammaire-interrogation-avec-et.pdf}} & 1er trimestre & 05/10 – 06/11/2026 & 91 \\ \hline
14 & Leçon 14 — Expression écrite : le paludisme \newline {\small\emph{Type : ecrit — fichier : ZH14-ecrit-le-paludisme.pdf}} & 1er trimestre & 05/10 – 06/11/2026 & 72 \\ \hline
15 & Leçon 15 — Expression orale : recyclage et paludisme \newline {\small\emph{Type : oral — fichier : ZH15-oral-recyclage-et-paludisme.pdf}} & 1er trimestre & 05/10 – 06/11/2026 & 105 \\ \hline
16 & Leçon 16 — Éléments socioculturels : maladies endémiques en Afrique et en Asie \newline {\small\emph{Type : socio — fichier : ZH16-socio-maladies-endemiques-en-afrique-et-en-asie.pdf}} & 1er trimestre & 05/10 – 06/11/2026 & 79 \\ \hline
\end{tabularx}\bigskip
\section*{Module 3 : Citoyenneté et ouverture au monde}
\textit{8 leçons, leçons 17 à 24, 690 pages.}\par\medskip
\begin{tabularx}{\linewidth}{|c|X|l|l|c|}\hline
\rowcolor{gray!20} N° & Titre de la leçon & Trimestre & Période & Pages \\ \hline
17 & Leçon 17 — Vocabulaire et compréhension de texte : collectivités locales \newline {\small\emph{Type : vocabulaire — fichier : ZH17-vocabulaire-collectivites-locales.pdf}} & 1er trimestre & 09/11 – 11/12/2026 & 81 \\ \hline
18 & Leçon 18 — Grammaire : 被字句 \newline {\small\emph{Type : grammaire — fichier : ZH18-grammaire.pdf}} & 1er trimestre & 09/11 – 11/12/2026 & 92 \\ \hline
19 & Leçon 19 — Expression écrite : offre de services aux usagers des administrations (collectivités locales) \newline {\small\emph{Type : ecrit — fichier : ZH19-ecrit-offre-de-services-aux-usagers-des-administra.pdf}} & 1er trimestre & 09/11 – 11/12/2026 & 95 \\ \hline
20 & Leçon 20 — Expression orale : collectivités locales \newline {\small\emph{Type : oral — fichier : ZH20-oral-collectivites-locales.pdf}} & 1er trimestre & 09/11 – 11/12/2026 & 87 \\ \hline
21 & Leçon 21 — Vocabulaire et compréhension de texte : démocratie \newline {\small\emph{Type : vocabulaire — fichier : ZH21-vocabulaire-democratie.pdf}} & 1er trimestre & 09/11 – 11/12/2026 & 79 \\ \hline
22 & Leçon 22 — Grammaire : phrase complexes \newline {\small\emph{Type : grammaire — fichier : ZH22-grammaire-phrase-complexes.pdf}} & 1er trimestre & 09/11 – 11/12/2026 & 88 \\ \hline
23 & Leçon 23 — Expression écrite : la démocratie \newline {\small\emph{Type : ecrit — fichier : ZH23-ecrit-la-democratie.pdf}} & 1er trimestre & 09/11 – 11/12/2026 & 88 \\ \hline
24 & Leçon 24 — Expression orale : démocratie \newline {\small\emph{Type : oral — fichier : ZH24-oral-democratie.pdf}} & 1er trimestre & 09/11 – 11/12/2026 & 80 \\ \hline
\end{tabularx}\bigskip
\section*{Module 4 : Activités économiques et monde du travail}
\textit{9 leçons, leçons 25 à 33, 781 pages.}\par\medskip
\begin{tabularx}{\linewidth}{|c|X|l|l|c|}\hline
\rowcolor{gray!20} N° & Titre de la leçon & Trimestre & Période & Pages \\ \hline
25 & Leçon 25 — Vocabulaire et compréhension de texte : gestion du budget \newline {\small\emph{Type : vocabulaire — fichier : ZH25-vocabulaire-gestion-du-budget.pdf}} & 2e trimestre & 04/01 – 19/02/2027 & 90 \\ \hline
26 & Leçon 26 — Grammaire : L’addition et soustraction en chinois 加法和减法 ; les différentes monnaies ; Les nombres ; 除了……以外，…… \newline {\small\emph{Type : grammaire — fichier : ZH26-grammaire-laddition-et-soustraction-en-chinois-les.pdf}} & 2e trimestre & 04/01 – 19/02/2027 & 70 \\ \hline
27 & Leçon 27 — Expression écrite : gestion de l'argent de poche et caractères complexes du budget \newline {\small\emph{Type : ecrit — fichier : ZH27-ecrit-gestion-de-l-argent-de-poche-et-caracteres-c.pdf}} & 2e trimestre & 04/01 – 19/02/2027 & 99 \\ \hline
28 & Leçon 28 — Expression orale : gestion du budget \newline {\small\emph{Type : oral — fichier : ZH28-oral-gestion-du-budget.pdf}} & 2e trimestre & 04/01 – 19/02/2027 & 113 \\ \hline
29 & Leçon 29 — Éléments socioculturels : différence entre la monnaie chinoise et la monnaie camerounaise \newline {\small\emph{Type : socio — fichier : ZH29-socio-difference-entre-la-monnaie-chinoise-et-la-m.pdf}} & 2e trimestre & 04/01 – 19/02/2027 & 91 \\ \hline
30 & Leçon 30 — Vocabulaire et compréhension de texte : publicité \newline {\small\emph{Type : vocabulaire — fichier : ZH30-vocabulaire-publicite.pdf}} & 2e trimestre & 04/01 – 19/02/2027 & 75 \\ \hline
31 & Leçon 31 — Grammaire : Les connecteurs 首先… 然后……其次……最后… ; 既…..又 ; 只有….才…. ; 只要……就….. \newline {\small\emph{Type : grammaire — fichier : ZH31-grammaire-les-connecteurs.pdf}} & 2e trimestre & 04/01 – 19/02/2027 & 79 \\ \hline
32 & Leçon 32 — Expression écrite : affiche publicitaire sur un produit (marketing) \newline {\small\emph{Type : ecrit — fichier : ZH32-ecrit-affiche-publicitaire-sur-un-produit-marketin.pdf}} & 2e trimestre & 04/01 – 19/02/2027 & 86 \\ \hline
33 & Leçon 33 — Éléments socioculturels : le marketing en Chine et au Cameroun \newline {\small\emph{Type : socio — fichier : ZH33-socio-le-marketing-en-chine-et-au-cameroun.pdf}} & 2e trimestre & 04/01 – 19/02/2027 & 78 \\ \hline
\end{tabularx}\bigskip
\section*{Module 5 : Mass médias et télécommunication}
\textit{7 leçons, leçons 34 à 40, 612 pages.}\par\medskip
\begin{tabularx}{\linewidth}{|c|X|l|l|c|}\hline
\rowcolor{gray!20} N° & Titre de la leçon & Trimestre & Période & Pages \\ \hline
34 & Leçon 34 — Vocabulaire et compréhension de texte : internet et les réseaux sociaux \newline {\small\emph{Type : vocabulaire — fichier : ZH34-vocabulaire-internet-et-les-reseaux-sociaux.pdf}} & 2e trimestre & 22/02 – 26/02/2027 & 98 \\ \hline
35 & Leçon 35 — Grammaire : 不管是…….或者是……. ; 是。。。。。。的 ; 比字句 \newline {\small\emph{Type : grammaire — fichier : ZH35-grammaire.pdf}} & 3e trimestre & 08/03 – 14/05/2027 & 86 \\ \hline
36 & Leçon 36 — Expression écrite : caractères complexes liés à la gestion des réseaux sociaux \newline {\small\emph{Type : ecrit — fichier : ZH36-ecrit-caracteres-complexes-lies-a-la-gestion-des-r.pdf}} & 3e trimestre & 08/03 – 14/05/2027 & 68 \\ \hline
37 & Leçon 37 — Éléments socioculturels : les outils TICs de marques chinoises et camerounaises \newline {\small\emph{Type : socio — fichier : ZH37-socio-les-outils-tics-de-marques-chinoises-et-came.pdf}} & 3e trimestre & 08/03 – 14/05/2027 & 87 \\ \hline
38 & Leçon 38 — Vocabulaire et compréhension de texte : rédaction d’un CV \newline {\small\emph{Type : vocabulaire — fichier : ZH38-vocabulaire-redaction-dun-cv.pdf}} & 3e trimestre & 08/03 – 14/05/2027 & 86 \\ \hline
39 & Leçon 39 — Grammaire : verbe modal 可以 ; 的 、 得、 地 \newline {\small\emph{Type : grammaire — fichier : ZH39-grammaire-verbe-modal.pdf}} & 3e trimestre & 08/03 – 14/05/2027 & 71 \\ \hline
40 & Leçon 40 — Expression écrite : caractères complexes liés à la rédaction d'un CV \newline {\small\emph{Type : ecrit — fichier : ZH40-ecrit-caracteres-complexes-lies-a-la-redaction-d-u.pdf}} & 3e trimestre & 08/03 – 14/05/2027 & 116 \\ \hline
\end{tabularx}\bigskip
\end{document}
